Un vaisseau s'éloigne de la Terre à $\num{\v{u1}}\,c$. Il lance une sonde vers l'avant, à $\num{\v{u2}}\,c$ par rapport au vaisseau.

\begin{enumerate}[label=\alph*)]
    \item Quelle est la vitesse de la sonde mesurée depuis la Terre ? Comparez avec le résultat de la relativité galiléenne.
    \item Si la sonde était plutôt lancée vers l'arrière (vers la Terre) à $\num{\v{u2}}\,c$ par rapport au vaisseau, quelle serait sa vitesse mesurée depuis la Terre, et dans quel sens irait-elle ?
    \item Deux vaisseaux identiques s'éloignent de la Terre en sens opposés, chacun à $\num{\v{u1}}\,c$ par rapport à la Terre. Quelle est la vitesse de l'un mesurée par l'autre ?
\end{enumerate}
